% Section: tautological-instances

\documentclass[../../../include/open-logic-section]{subfiles}

\begin{document}

\olfileid[fr]{nml}{syn}{tau}

\olsection{Instances tautologiques}

\begin{explain}
Une formule sans opérateur modal est une tautologie si elle est vraie
pour toute attribution de valeurs de vérité. Toute tautologie est
donc vraie en chaque monde de chaque modèle. Mais la notion de
tautologie n'est pas définie pour les !!{formula}s contenant
$\Box$ ou~$\Diamond$. Par exemple, $\Box p \lor \lnot \Box p$,
qui est une instance du principe du tiers exclu, est-elle valide ?
La notion d'\emph{instance tautologique} permet de répondre :
il s'agit d'une !!{formula} obtenue par substitution à partir d'une
tautologie sans opérateur modal. Toute instance tautologique est
valide, ce qui n'a rien de surprenant, mais demande une preuve.
\end{explain}

\begin{defn}
  Une !!{formula} modale~$!B$ est une \emph{instance tautologique}
  si et seulement s'il existe une tautologie sans opérateur
  modal~$!A$, dont les !!{propositional variable}s sont $p_1$,
  \dots,~$p_n$, et des !!{formula}s $!D_1$, \dots,~$!D_n$ telles
  que $!B \ident \SSubst{!A}{\subst{!D_1}{p_1},
    \dots, \subst{!D_n}{p_n}}$.
\end{defn}

\begin{lem}\ollabel{lem:valid-taut}
  Soit $!A$ une !!{formula} sans opérateur modal dont les
  !!{propositional variable}s sont $p_1$, \dots, $p_n$, et
  soient $!D_1$, \dots, $!D_n$ des !!{formula}s modales.
  Pour toute attribution $\pAssign{v}$, tout modèle
  $\mModel{M} = \tuple{W, R, V}$ et tout $w \in W$ tels que
  $\pAssign{v}(p_i) = \True$ si et seulement si
  $\mSat{M}{!D_i}[w]$, nous avons $\pSat{v}{!A}$ si et seulement si
  $\mSat{M}{\SSubst{!A}{\subst{!D_1}{p_1}, \dots,
      \subst{!D_n}{p_n}}}[w]$.
\end{lem}

\begin{proof}
  Par induction sur~$!A$.
  \begin{enumerate}
  \tagitem{prvFalse}{\indcase{!A}{\lfalse}{On a à la fois
      $\pSat/{v}{\lfalse}$ et $\mSat/{M}{\lfalse}[w]$}.}{}
  \tagitem{prvTrue}{\indcase{!A}{\ltrue}{On a à la fois
      $\pSat{v}{\ltrue}$ et $\mSat{M}{\ltrue}[w]$}.}{}
  \item \indcase{!A}{p_i}{%
    \begin{align*}
      \pSat{v}{p_i} \Leftrightarrow {} & \pAssign{v}(p_i) = \True \\
      & \qquad \text{par définition de $\pSat{v}{p_i}$}\\
      \Leftrightarrow {} & \mSat{M}{!D_i}[w] \\
      &\qquad \text{par hypothèse}\\
      \Leftrightarrow {} & \mSat{M}{\SSubst{p_i}{\subst{!D_1}{p_1}, \dots,
          \subst{!D_n}{p_n}}}[w]\\
      &\qquad \text{car $\SSubst{p_i}{\subst{!D_1}{p_1}, \dots,
          \subst{!D_n}{p_n}} \ident !D_i$}.
      \end{align*}}
  \tagitem{prvNot}{\indcase{!A}{\lnot !B}{%
      \begin{align*}
        \pSat{v}{\lnot !B} \Leftrightarrow {} & \pSat/{v}{!B}\\
        &\qquad \text{par définition de $\pSat{v}{}$};\\
        \Leftrightarrow {} & \mSat/{M}{\SSubst{!B}{\subst{!D_1}{p_1}, \dots,
          \subst{!D_n}{p_n}}}[w]\\
        &\qquad \text{par hypothèse d'induction}\\
        \Leftrightarrow {} &
        \mSat{M}{\SSubst{\lnot !B}{\subst{!D_1}{p_1}, \dots,
          \subst{!D_n}{p_n}}}[w]\\
        &\qquad \text{par définition de $\mSat{M}{}[w]$}.
      \end{align*}}}{}
  \tagitem{prvAnd}{\indcase{!A}{(!B \land !C)}{%
      \begin{align*}
        \pSat{v}{!B \land !C} \Leftrightarrow {} &
        \pSat{v}{!B} \text{ et } \pSat{v}{!C}\\
        &\qquad \text{par définition de $\pSat{v}{}$}\\
        \Leftrightarrow {} &
        \mSat{M}{\SSubst{!B}{\subst{!D_1}{p_1}, \dots,
            \subst{!D_n}{p_n}}}[w] \text{ et } \\
        & \mSat{M}{\SSubst{!C}{\subst{!D_1}{p_1}, \dots,
          \subst{!D_n}{p_n}}}[w]\\
        &\qquad \text{par hypothèse d'induction}\\
        \Leftrightarrow{} &
        \mSat{M}{\SSubst{(!B \land !C)}{\subst{!D_1}{p_1}, \dots,
          \subst{!D_n}{p_n}}}[w]\\
        &\qquad \text{par définition de $\mSat{M}{}[w]$}.
      \end{align*}}}{}
  \tagitem{prvOr}{\indcase{!A}{(!B \lor !C)}{%
      \begin{align*}
        \pSat{v}{!B \lor !C} \Leftrightarrow{} &
        \pSat{v}{!B} \text{ ou } \pSat{v}{!C}\\
        &\qquad \text{par définition de $\pSat{v}{}$};\\
        \Leftrightarrow{} & \mSat{M}{\SSubst{!B}{\subst{!D_1}{p_1}, \dots,
            \subst{!D_n}{p_n}}}[w] \text{ ou }\\
        & \mSat{M}{\SSubst{!C}{\subst{!D_1}{p_1}, \dots,
          \subst{!D_n}{p_n}}}[w]\\
        &\qquad \text{par hypothèse d'induction}\\
        \Leftrightarrow{} &
        \mSat{M}{\SSubst{(!B \lor !C)}{\subst{!D_1}{p_1}, \dots,
          \subst{!D_n}{p_n}}}[w]\\
        &\qquad \text{par définition de $\mSat{M}{}[w]$}.
      \end{align*}}}{}
  \tagitem{prvIf}{\indcase{!A}{(!B \lif !C)}{%
      \begin{align*}
        \pSat{v}{!B \lif !C} \Leftrightarrow{} &
        \pSat/{v}{!B} \text{ ou } \pSat{v}{!C}\\
        &\qquad \text{par définition de $\pSat{v}{}$}\\
        \Leftrightarrow{} &
        \mSat/{M}{\SSubst{!B}{\subst{!D_1}{p_1}, \dots,
            \subst{!D_n}{p_n}}}[w] \text{ ou }\\
        & \mSat{M}{\SSubst{!C}{\subst{!D_1}{p_1}, \dots,
          \subst{!D_n}{p_n}}}[w]\\
        &\qquad \text{par hypothèse d'induction}\\
        \Leftrightarrow{} &
        \mSat{M}{\SSubst{(!B \lif !C)}{\subst{!D_1}{p_1}, \dots,
          \subst{!D_n}{p_n}}}[w]\\
        &\qquad \text{par définition de $\mSat{M}{}[w]$}.
      \end{align*}}}{}
  \tagitem{prvIff}{\indcase{!A}{(!B \liff !C)}{%
      \begin{align*}
        \pSat{v}{!B \liff !C} \Leftrightarrow {} &
        \text{soit } \pSat{v}{!B} \text{ et } \pSat{v}{!C}\\
        & \text{soit } \pSat/{v}{!B} \text{ et } \pSat/{v}{!C}\\
        &\qquad \text{par définition de $\pSat{v}{}$}\\
        \Leftrightarrow {} &
        \text{soit } \mSat{M}{\SSubst{!B}{\subst{!D_1}{p_1}, \dots,
            \subst{!D_n}{p_n}}}[w] \text{ et }\\
        & \qquad \mSat{M}{\SSubst{!C}{\subst{!D_1}{p_1}, \dots,
          \subst{!D_n}{p_n}}}[w]\\ 
        &  \text{soit } \mSat/{M}{\SSubst{!B}{\subst{!D_1}{p_1}, \dots,
            \subst{!D_n}{p_n}}}[w] \text{ et }\\
        & \qquad \mSat/{M}{\SSubst{!C}{\subst{!D_1}{p_1}, \dots,
          \subst{!D_n}{p_n}}}[w]\\
        &\qquad\text{par hypothèse d'induction}\\
        \Leftrightarrow {} &
        \mSat{M}{\SSubst{(!B \liff !C)}{\subst{!D_1}{p_1}, \dots,
          \subst{!D_n}{p_n}}}[w]\\
        &\qquad \text{par définition de $\mSat{M}{}[w]$}.
      \end{align*}}}{}
  \end{enumerate}
  \end{proof}

\begin{prop}\ollabel{prop:valid-taut}
  Toutes les instances tautologiques sont valides.
\end{prop}

\begin{proof}
  Raisonnons par contraposition. Supposons qu'il existe $!A$ tel que
  $\mSat/{M}{\SSubst{!A}{\subst{!D_1}{p_1}, \dots,
  \subst{!D_n}{p_n}}}[w]$, pour un certain modèle $\mModel{M}$ et
  un certain monde~$w$. Définissons une attribution $\pAssign{v}$
  telle que $\pAssign{v}(p_i) = \True$ si et seulement si
  $\mSat{M}{!D_i}[w]$ (et $\pAssign{v}$ attribue des valeurs
  arbitraires aux $q \notin \{p_1, \dots, p_n \}$). D'après le
  \olref{lem:valid-taut}, $\pSat/{v}{!A}$; donc $!A$ n'est pas une
  tautologie.
\end{proof}
\begin{editorial}Dans la preuve du lemme, le cas de la négation portait dans la source le tag de la constante fausse ; il est replacé sous la négation. Sa dernière justification renvoyait à la satisfaction propositionnelle, au lieu de la satisfaction dans un modèle en un monde. Enfin, la première ligne du cas du biconditionnel affichait le conditionnel. Ces trois indications sont corrigées ; les étapes restantes et la conclusion sont conservées.\end{editorial}
\end{document}
